\subsection{Analyse des résultats}
\label{sec:DMAanalyse}
Grâce au traitement des données sur la Fig.~\ref{fig:DMA_resultats}, nous pouvons obtenir, en couplant les extrema à chaque fréquence d’oscillation imposée à la masse, le tableau suivant :

\\
\begin{table}[!ht]
\centering
\footnotesize
\setlength{\tabcolsep}{4pt}
\renewcommand{\arraystretch}{0.9}
\begin{tabular}{|c|c|c|}
\hline
	\textbf{Frequence [Hz]} 
	& \textbf{Peak [$^\circ$]}
	& \textbf{incertitude [$^\circ$]}
\\\hline
	1.000 & 140.5 & 1.606 \\
	2.000 & 143.5 & 1.594 \\
	5.000 & 148.6 & 1.601 \\
	10.000 & 150.5 & 1.616 \\
	20.000 & 155.7 & 1.518 \\
\hline
\textbf{Moyenne}
& \textbf{n/a}
& \textbf{1.587} \\
\hline
\end{tabular}
	\caption{Valeurs de temperatures auxquelles E (module d'Young) est maximal}
\label{tab:peaks}
\end{table}
On peut donc extraire tous les extrema locaux de la Fig. \ref{fig:DMA_resultats} que l’on peut retrouver dans le tableau \ref{tab:peaks}. Grâce à la relation Eq. \ref{eq:arrhenius} que l’on linéarise tel que Eq. \ref{eq:arrheniuslin}, nous obtenons le graphique d’Arrhenius suivant :

\begin{figure}[hct]
        \centering
        \includegraphics[scale=0.7]{my_arrenius.png}
        \label{fig:arrenius}
        \caption{Arrhenius graph}
\end{figure}
que l’on peut comparer avec celui obtenu grâce au logiciel du fournisseur (voir Ref. \ref{fig:n_arrenius}). Le graphique trace le logarithme naturel de la fréquence en fonction de l'inverse de la température. Cela nous permet, comme nous le mentionnons en Sec. \ref{sec:theorie}, d’extraire l’énergie d’activation $E_a$ associée à la transition vitreuse de notre échantillon de PMMA, soit :
\[
	E_a &= 295.730 \pm 18.959 ~[kJ/mol].
\]
